\subsection{Preferences and a restricted execution interface}
A concrete preference interpretation of expected participation is quasilinear, risk-neutral continuation utility $A_j-\E(Y+C)$ with outside value $O_j$, so that $b_j=A_j-O_j$ after normalization. Alternatively the expectation bound can be an explicitly agreed actuarial budget. Neither interpretation follows from a finite communication alphabet. A customer with a hard realization-level liquidity limit instead requires $Y+C\leq b_j$. The intermediate institution retains the actuarial cap and separately authorizes at most $\delta$ realization overrun. Its intermediate tier is committed before the draw so that the disclosed lottery and its ceiling can be checked together. Allowing a fresh, costless rejection after a high draw changes the feasible problem; the expected-participation comparison assumes that option is absent, not that every real renewal contract can eliminate it.

The downstream restriction is an interface restriction, not an inability to store a model. An upstream desk keeps the branch record; a downstream gateway has a read-only map from a small session token to tested execution levels and no access to the upstream record or a branch-indexed customer identifier. A relevant documented component analogy is operational-technology allowlisting and least functionality: NIST SP~800-82 Rev.~3 discusses restricting executable applications, testing policy configurations, disabling nonessential interfaces, and restricting writable set points \citep[Sections~5.2.5.1, 6.2.4.1 and control CM-7]{StoufferEtAl2023}. These are precedents for a constrained execution boundary, not evidence that a particular provider uses our entire renewal or lottery protocol.

Under this design, opening an additional level can incur a one-time cost of engineering, test fixtures, authorization, and configuration assurance. This is the modeled source of $\rho(c)$: it is paid if the level is enabled, independently of the frequency with which its token is used. Per-use charges and entropy-based communication charges would be different objective terms. An exogenous catalog represents levels already eligible for certification; a numerical catalog chosen by the operator instead calls for the approximation theorem in Section~\ref{sec:grid}. The model does not assert that every read-only contract must be inaccessible downstream; it quantifies the consequence when that access is deliberately excluded. No empirical adoption, calibration, or enforceability conclusion is inferred from the analogy.
